\documentclass[11pt]{article}
\usepackage[margin=0.85in]{geometry}
\usepackage{amsmath,amssymb,amsfonts}
\usepackage{booktabs,array,xcolor}
\usepackage{hyperref}
\hypersetup{colorlinks=true,linkcolor=blue,urlcolor=blue}
\newcommand{\Lam}{\Lambda}
\title{\textbf{Why LGM-SSP cannot beat LGM}\\[2pt]
\large A mechanistic analysis of the linear latent-Hawkes vs.\ the rate-pinned soft-max model}
\author{}\date{}
\begin{document}\maketitle\vspace{-3em}

\section{Headline}
\textbf{LGM-SSP traded away expressivity for a calibration property LGM already had.}
LGM is \emph{already} exactly rate-pinned (scalar $\mu_0=R_{\mathrm{target}}(1-n)$) \emph{and} keeps a
deep nonlinear mark head. LGM-SSP, to put everything in one linear latent-Hawkes field, made
\textbf{both the rate and the marks linear} --- and gained \emph{no} calibration advantage in return
(LGM was already calibrated). So it is a near-pure expressivity loss. The two models are
\[
\textbf{LGM: }\ \lambda_k=\underbrace{\Lam(t)}_{\text{linear scalar Hawkes, pinned}}\!\cdot\,
\underbrace{\operatorname{softmax}(z_k)}_{\text{DEEP NONLINEAR per-type net}}
\qquad\qquad
\textbf{LGM-SSP: }\ \lambda_k=\underbrace{\nu_k+(Wx)_k}_{\text{LINEAR readout of one latent field}} .
\]

\section{Measured gaps (cbse-BTC days 1--2, seq64)}
\begin{center}\small\renewcommand{\arraystretch}{1.2}
\begin{tabular}{lcccc}
\toprule
model & ACC $\uparrow$ & PPL $\downarrow$ & SIM $\downarrow$ & clustering $\mathtt{clus\_re}\downarrow$\\
\midrule
LGM \texttt{L64\_r92\_M8} (nonlinear, best pred.) & \textbf{0.4322} & \textbf{8.80} & 0.357 & \textbf{0.12}\\
LGM \texttt{L128\_r92\_M4\_vfb} (best sim.) & 0.350 & 11.2 & \textbf{0.14} & 0.04\\
LGM-SSP \texttt{P256\_n99} (best tuned) & 0.3846 & 10.49 & 0.822 & 0.94\\
LGM-SSP \texttt{P128\_n97} (best-sim tuned) & 0.358 & 10.89 & 0.748 & 0.93\\
\bottomrule
\end{tabular}
\end{center}

\section{Four mechanisms}
\paragraph{1. Marks: deep nonlinear soft-max vs.\ a ratio of linear functions $\Rightarrow$ the prediction gap ($0.432\to0.385$).}
Next-mark accuracy is governed by $p(k\mid t)$. LGM uses $p(k)=\operatorname{softmax}(z_k(h))$ with
$z_k$ from a \emph{deep per-type recurrent network} --- it can learn rich history$\to$type maps. LGM-SSP
uses $p(k)=\lambda_k/\sum_j\lambda_j=(\nu_k+Wx_k)/\sum_j(\nu_j+Wx_j)$, a \textbf{ratio of linear
functions of one shared $x$}, with no nonlinear capacity and tied to the rate field. A linear logit
underfits a 62-way conditional categorical with complex LOB dependencies --- the dominant driver of
the accuracy/perplexity gap.

\paragraph{2. Rate: nonlinear vol-feedback vs.\ linear self-excitation $\Rightarrow$ the clustering gap ($\mathtt{clus\_re}\ 0.93\to0.12$).}
Volatility clustering (the $|r|$-ACF) needs the rate to \emph{amplify during directional bursts} --- a
multiplicative, state-dependent effect. LGM($+$vfb) adds a \textbf{nonlinear} QHawkes term
$b(R^2-c_0)$ so the rate spikes during signed-flow runs (fat tails $+$ clustering). LGM-SSP's
$\Lam=\sum_k(\nu_k+Wx_k)$ is \textbf{purely linear}: a linear Hawkes gives only exponential,
branching-bounded clustering and cannot match real conditional-variance clustering. Hence
$\mathtt{clus\_re}$ stays $\approx0.93$ \emph{across all $P$ and $n_{\mathrm{cap}}$} --- structural, not
a hyperparameter.

\paragraph{3. Entanglement vs.\ decoupling.}
LGM \emph{separates} the jobs: the scalar rate handles \emph{when/how-often} (and is calibrated); the
free nonlinear net handles \emph{which type}. Each uses its full capacity for its task. LGM-SSP
\emph{fuses} rate and marks into one linear field, so the mark head is constrained by the rate's
structure and cannot specialize. Decoupling is precisely why LGM's marks can be expressive \emph{and}
its rate stays calibrated at the same time.

\paragraph{4. The $n<1$ constraint caps the only clustering LGM-SSP has.}
LGM-SSP's clustering comes entirely from linear self-excitation, whose strength \emph{is} the branching
ratio $n$. But stability/calibration require $n<1$, so the architecture \textbf{caps the very mechanism}
it depends on. Empirically: $n_{\mathrm{cap}}=0.90\Rightarrow$ collapse (ACC $0.17$, too damped); even
$n_{\mathrm{cap}}=0.99$ (near-critical) cannot cluster enough. LGM's clustering comes from the
\emph{nonlinear} vfb, which is \emph{not} subject to the linear-branching cap --- it can cluster hard
while staying calibrated.

\section{The tuning sweep isolates the cause}
\begin{center}\small\renewcommand{\arraystretch}{1.2}
\begin{tabular}{>{\raggedright\arraybackslash}p{5.2cm} >{\raggedright\arraybackslash}p{5.2cm} >{\raggedright\arraybackslash}p{4.3cm}}
\toprule
lever changed & expectation \emph{if} it were the bottleneck & observed\\
\midrule
$n_{\mathrm{cap}}:0.90\to0.99$ & clustering improves & ACC recovers, \textbf{clustering flat $\approx0.93$}\\
$P:128\to256$ (more latent modes) & richer kernel $\Rightarrow$ clustering improves & \textbf{prediction $+0.01$, clustering flat}\\
\bottomrule
\end{tabular}
\end{center}
Neither branching nor capacity moves clustering $\Rightarrow$ the limit is \textbf{linearity}, not
hyperparameters.

\section{The deeper point and conclusion}
The one thing LGM-SSP uniquely adds over LGM is \textbf{cross-excitation} (mark $j$ raising mark $k$'s
rate). But because it is \emph{linear}, that cross-excitation is weak, and LGM already captures
cross-mark structure inside its nonlinear mark soft-max --- so LGM-SSP's single novel ingredient does
not pay for the expressivity it gave up.

\medskip
\noindent\textbf{Conclusion.} LGM-SSP cannot beat LGM because it removed the two nonlinearities that
matter --- the deep mark head (drives prediction) and the vol-feedback rate (drives clustering) --- in
exchange for a calibration guarantee LGM \emph{already} provides. \textbf{Exact calibration does not
require linearity}: LGM gets the exact mean pin with a \emph{scalar} rate while keeping nonlinear marks;
LGM-SSP linearized everything and paid for it with no calibration upside. To close the gap one must
re-introduce nonlinearity (a mean-preserving multiplicative modulation, or a nonlinear mark head), at
which point the model converges back toward LGM / S2P2.
\end{document}
